\rhead{CY202: Cryptography \& Network Security}
\phantomsection 
\section*{Cryptography \& Network Security (CY202)}
\label{major:CS_2}

\noindent \small{\hyperref[main:scheme]{$\leftarrow$ Back to Scheme}} \vspace{0.5cm}

\noindent \textbf{Credits:} 3 (2-1-0)

\subsection*{Theory}
\noindent\textbf{Unit 1} \\
\textit{Cryptography Fundamentals and Classical Techniques:} Security services (confidentiality, integrity, authentication, non-repudiation), Classical ciphers (Caesar, Vigenère, Playfair, Hill), Cryptanalysis techniques (frequency analysis, known-plaintext attacks), Block vs. stream ciphers, Shannon's principles of confusion and diffusion, Modes of block cipher operation (ECB, CBC, CFB, OFB, CTR), Initialization vectors and nonce usage.

\vspace{0.3cm}

\noindent\textbf{Unit 2} \\
\textit{Symmetric Key Cryptography:} Data Encryption Standard (DES) algorithm and cryptanalysis (differential, linear), Advanced Encryption Standard (AES) - Rijndael structure, S-boxes, key expansion, rounds structure, Block cipher modes security analysis, Key management challenges in symmetric cryptography, Key derivation functions and key stretching, Practical symmetric cipher selection criteria.

\vspace{0.3cm}

\noindent\textbf{Unit 3} \\
\textit{Number Theory and Public-Key Cryptography:} Modular arithmetic, prime generation, Fermat/Euler theorems, Greatest Common Divisor algorithms, RSA cryptosystem (key generation, encryption/decryption, security proofs), Diffie-Hellman key exchange and man-in-the-middle attacks, Digital Signature Algorithm (DSA), Elliptic Curve Cryptography (ECC) fundamentals and advantages over RSA.

\vspace{0.3cm}

\noindent\textbf{Unit 4} \\
\textit{Hash Functions, MACs, and Authentication:} Cryptographic hash function properties (collision resistance, preimage resistance), MD5 weaknesses and SHA-family algorithms, HMAC construction and security proofs, Message authentication codes (MACs), Digital signatures and non-repudiation, PKI concepts (certificates, CAs, CRLs), X.509 certificate format and validation.

\vspace{0.3cm}

\noindent\textbf{Unit 5} \\
\textit{Network Security Protocols and Applications:} Transport Layer Security (TLS/SSL) handshake and cipher suites, IPsec architecture (AH, ESP, IKE), VPN protocols and tunnel modes, Email security (PGP, S/MIME), Kerberos authentication and ticket granting, Common network attacks (MITM, replay, session hijacking), Secure protocol design principles and cryptographic agility.

\newpage

\subsection*{Practical}
\noindent\textbf{Lab 1: Classical Cipher Implementation} \\
Focuses on implementing and cryptanalyzing classical ciphers (Caesar, Vigenère), demonstrating frequency analysis attacks and basic cryptanalysis techniques.

\vspace{0.2cm}

\noindent\textbf{Lab 2: Block Cipher Modes Analysis} \\
Focuses on implementing ECB/CBC/CTR modes, demonstrating padding oracle attacks on ECB and malleability issues in CBC mode.

\vspace{0.2cm}

\noindent\textbf{Lab 3: AES Implementation Study} \\
Focuses on studying AES encryption/decryption process, key expansion, and analyzing performance characteristics across different key sizes.

\vspace{0.2cm}

\noindent\textbf{Lab 4: RSA Key Generation and Attacks} \\
Focuses on RSA implementation, textbook vs. padded RSA, demonstrating small exponent attacks and common modulus attacks.

\vspace{0.2cm}

\noindent\textbf{Lab 5: Diffie-Hellman Key Exchange} \\
Focuses on implementing Diffie-Hellman key exchange, demonstrating man-in-the-middle attacks and ephemeral key usage.

\vspace{0.2cm}

\noindent\textbf{Lab 6: Hash Function Collision Attacks} \\
Focuses on demonstrating MD5 collision attacks and analyzing SHA-256 preimage resistance through birthday paradox experiments.

\vspace{0.2cm}

\noindent\textbf{Lab 7: HMAC and Message Authentication} \\
Focuses on implementing HMAC construction and verifying security properties against length extension attacks.

\vspace{0.2cm}

\noindent\textbf{Lab 8: Digital Signature Verification} \\
Focuses on creating/verifying digital signatures using RSA/ECDSA, implementing X.509 certificate chain validation.

\vspace{0.2cm}

\noindent\textbf{Lab 9: TLS Handshake Analysis} \\
Focuses on capturing and analyzing TLS handshakes using Wireshark, identifying cipher suite negotiation and certificate validation.

\vspace{0.2cm}

\noindent\textbf{Lab 10: Network Security Protocol Implementation} \\
Focuses on implementing simplified TLS/IPsec protocol components, demonstrating secure key exchange and data protection.

\vspace{0.2cm}

\newpage
\rhead{Curriculum Schema}
